%-------------------------
% Resume in Latex
% Author : Levon Khorasandzhian
% Based off of: https://github.com/sb2nov/resume
% License : MIT
%------------------------

\documentclass[letterpaper,11pt]{article}

\usepackage{latexsym}
\usepackage[empty]{fullpage}
\usepackage{titlesec}
\usepackage{marvosym}
\usepackage[usenames,dvipsnames]{color}
\usepackage{verbatim}
\usepackage{enumitem}
\usepackage[hidelinks]{hyperref}
\usepackage{fancyhdr}
\usepackage[T2A]{fontenc}
\usepackage[utf8]{inputenc}
\usepackage[english, russian]{babel}
\usepackage{tabularx}
\setlength{\footskip}{4.08003pt}
\input{glyphtounicode}


%----------FONT OPTIONS----------
% sans-serif
% \usepackage[sfdefault]{FiraSans}
% \usepackage[sfdefault]{roboto}
% \usepackage[sfdefault]{noto-sans}
% \usepackage[default]{sourcesanspro}

% serif
% \usepackage{CormorantGaramond}
% \usepackage{charter}


\pagestyle{fancy}
\fancyhf{} % clear all header and footer fields
\fancyfoot{}
\renewcommand{\headrulewidth}{0pt}
\renewcommand{\footrulewidth}{0pt}

% Adjust margins
\addtolength{\oddsidemargin}{-0.5in}
\addtolength{\evensidemargin}{-0.5in}
\addtolength{\textwidth}{1in}
\addtolength{\topmargin}{-.5in}
\addtolength{\textheight}{1.0in}

\urlstyle{same}

\raggedbottom
\raggedright
\setlength{\tabcolsep}{0in}

% Sections formatting
\titleformat{\section}{
  \vspace{-4pt}\scshape\raggedright\large
}{}{0em}{}[\color{black}\titlerule \vspace{-5pt}]

% Ensure that generate pdf is machine readable/ATS parsable
\pdfgentounicode=1

%-------------------------
% Custom commands
\newcommand{\resumeItem}[1]{
  \item\small{
    {#1 \vspace{-4pt}}
  }
}

\newcommand{\resumeSubheading}[4]{
  \vspace{-2pt}\item
    \begin{tabular*}{0.97\textwidth}[t]{l@{\extracolsep{\fill}}r}
      \textbf{#1} & #2 \\
      \textit{\small#3} & \textit{\small #4} \\
    \end{tabular*}\vspace{-7pt}
}

\newcommand{\resumeSubSubheading}[2]{
    \item
    \begin{tabular*}{0.97\textwidth}{l@{\extracolsep{\fill}}r}
      \textit{\small#1} & \textit{\small #2} \\
    \end{tabular*}\vspace{-7pt}
}

\newcommand{\resumeProjectHeading}[2]{
    \item
    \begin{tabular*}{0.97\textwidth}{l@{\extracolsep{\fill}}r}
      \small#1 & #2 \\
    \end{tabular*}\vspace{-7pt}
}

\newcommand{\resumeSubItem}[1]{\resumeItem{#1}\vspace{-4pt}}

\renewcommand\labelitemii{$\vcenter{\hbox{\tiny$\bullet$}}$}

\newcommand{\resumeSubHeadingListStart}{\begin{itemize}[leftmargin=0.15in, label={}]}
\newcommand{\resumeSubHeadingListEnd}{\end{itemize}}
\newcommand{\resumeItemListStart}{\begin{itemize}}
\newcommand{\resumeItemListEnd}{\end{itemize}\vspace{-5pt}}

%-------------------------------------------
%%%%%%  RESUME STARTS HERE  %%%%%%%%%%%%%%%%%%%%%%%%%%%%


\begin{document}

%----------HEADING----------
% \begin{tabular*}{\textwidth}{l@{\extracolsep{\fill}}r}
%   \textbf{\href{http://sourabhbajaj.com/}{\Large Sourabh Bajaj}} & Email : \href{mailto:sourabh@sourabhbajaj.com}{sourabh@sourabhbajaj.com}\\
%   \href{http://sourabhbajaj.com/}{http://www.sourabhbajaj.com} & Mobile : +1-123-456-7890 \\
% \end{tabular*}

\begin{center}
    \textbf{\Huge \scshape Левон Хорасанджян} \\ \vspace{1pt}
    \href{mailto:levon.kh228@gmail.com}{\underline{levon.kh228@gmail.com}} $|$ 
    \href{https://t.me/yungflaeva}{\underline{t.me/yungflaeva}} $|$
    \href{https://github.com/lkhorasandzhian}{\underline{github.com/lkhorasandzhian}}
\end{center}


%-----------EDUCATION-----------
\section{Образование}
  \resumeSubHeadingListStart
    \resumeSubheading
      {Высшая Школа Экономики}{Москва, Россия}
      {Бакалавриат Факультета Компьютерных Наук, Программная Инженерия}{Август 2021 -- Июль 2025}
    \resumeSubSubheading
      {Магистратура Факультета Компьютерных Наук, Инженерия Данных}{Октябрь 2025 -- Июль 2027}
      \resumeItem
    {\textbf{Основные курсы (GPA 8.05/10)}:
    \begin{enumerate}
        \item Алгоритмы и структуры данных
        \item Архитектура программных систем
        \item Программирование на C\#
        \item Конструирование программного обеспечения на Java
        \item Распределенные базы данных и хранилища
        \item Инструменты промышленной разработки
        \item Контроль качества программного обеспечения и тестирование
    \end{enumerate}
    }
  \resumeSubHeadingListEnd

%-----------PROJECTS-----------
\section{Проекты}
    \resumeSubHeadingListStart
          \resumeProjectHeading
          {\href{https://github.com/lkhorasandzhian/ylab-university}{\underline{\textbf{Product Catalog Service}}} $|$ \emph{Java, Spring, PostgreSQL, Docker, Swagger, AOP, Web MVC}}{Ноябрь 2025}
          \resumeItemListStart
              \resumeItem{REST API для продуктового маркетплейса с авторизацией, ролями и разграничением доступа}
              \resumeItem{Хранение данных с помощью реляционной СУБД PostgreSQL с контейнеризацией через Docker}
              \resumeItem{Документация с использованием Swagger и Javadoc}
              \resumeItem{Сквозная логика (логирование, валидация, контроль доступа) с помощью Spring AOP}
              \resumeItem{Развитие образовательного проекта в 5 итерациях (Git-ветки): от консольного приложения и сервлетов до полноценного Spring MVC сервиса}
          \resumeItemListEnd
          
          \resumeProjectHeading
          {\href{https://github.com/lkhorasandzhian/aton-web-api}{\underline{\textbf{Aton Web API}}} $|$ \emph{C\#, ASP.NET Core 8.0}}{Август, 2024}
          \resumeItemListStart
            \resumeItem{Web API сервис на .NET Core 8.0, реализующий CRUD методы над сущностью Users}
            \resumeItem{Доступ осуществляется через интерфейс Swagger}
            \resumeItem{Для выполнения запросов требуется авторизация}
          \resumeItemListEnd
          
           \resumeProjectHeading
            {\href{https://github.com/lkhorasandzhian/restaurant-order-processing-system}{\underline{\textbf{Restaurant Order Processing System}}} $|$ \emph{Java, Spring, H2, JavaDoc, Postman}}{Май 2023}
            \resumeItemListStart
              \resumeItem{Разработка Backend-системы обработки заказов ресторана на основе микросервисной архитектуры}
              \resumeItem{Реализация двух микросервисов: авторизации пользователей и управления заказами}
              \resumeItem{Проектирование REST API для регистрации, аутентификации, управления ролями и просмотра меню}
              \resumeItem{Организация хранения данных в in-memory СУБД H2}
              \resumeItem{Подготовка документации кода с использованием JavaDoc и коллекции Postman для тестирования API}
          \resumeItemListEnd
          
          \resumeProjectHeading
          {\href{https://github.com/lkhorasandzhian/notepad-shell}{\underline{\textbf{Notepad Shell}}} $|$ \emph{C\#, Windows Forms}}{Март, 2023}
          \resumeItemListStart
            \resumeItem{Desktop-приложение текстовый блокнот для работы с документами}
            \resumeItem{Реализация операций с текстовыми файлами и настройка цветовой темы интерфейса}
            \resumeItem{Создание пользовательского интерфейса с помощью Windows Forms}
          \resumeItemListEnd

          \resumeProjectHeading
          {\href{https://github.com/lkhorasandzhian/valid_sudoku}{\underline{\textbf{Valid Sudoku}}} $|$ \emph{Dart, Flutter}}{Май, 2023}
          \resumeItemListStart
            \resumeItem{Кроссплатформенное Desktop-приложение для Windows/MacOS/Linux}
            \resumeItem{Автогенерация задач 3 видов сложности и последующая валидация решения пользователя}
            \resumeItem{Загрузка готовой задачи, валидация на корректность правилам игры судоку, а также проверка наличия единственного решения}
            \resumeItem{Реализация рекурсивного backtracking-алгоритма для проверки решений}
          \resumeItemListEnd
    \resumeSubHeadingListEnd


%---------competitions-------------

%
%-----------PROGRAMMING SKILLS-----------
\section{Технические навыки}
 \begin{itemize}[leftmargin=0.15in, label={}]
    \small{\item{
     \textbf{Языки}{: Java, C\#} \\
     \textbf{Фреймворки}{: \textbf{Spring} (Core, Web MVC, REST, Security, AOP) · \textbf{.NET} (ASP.NET Core, Web API, EF Core, LINQ, DI, Windows Forms)} \\
     \textbf{СУБД}{: PostgreSQL, MySQL, MongoDB} \\
     \textbf{Инструменты разработчика}{: Git, Docker, Postman, WSL, VirtualBox, Visual Paradigm} \\
     \textbf{Тестирование}{: JUnit 5, xUnit, Mockito, Testcontainers}
    }}
 \end{itemize}

%------------LANGUAGES
\section{Языки}
 \begin{itemize}[leftmargin=0.15in, label={}]
    \small{\item{
     \textbf{Английский}{: B2} \\
     \textbf{Русский}{: Native}
    }}
 \end{itemize}

%-------------------------------------------
\end{document}
